% 1. บทนำ -- ย่อจากรายงานบทที่ 1 (../report/chapter/01-introduction.tex)
\section{บทนำ}

นักศึกษาใช้โซเชียลมีเดียทุกวัน ทั้งเพื่อการเรียน การติดตามข่าว และการติดต่อกับเพื่อน
หากใช้มากเกินไปในช่วงที่นอนไม่พอและมีความกดดันจากการเรียนอยู่แล้ว สุขภาพจิตอาจแย่ลงโดยไม่รู้ตัว
นอกจากเวลาที่ใช้แล้ว ประเภทเนื้อหาที่เสพก็อาจมีผลด้วย
การไล่อ่านข่าวเชิงลบ เช่น การเมือง อาชญากรรม หรือภัยพิบัติ ต่อเนื่องเป็นเวลานาน
เรียกว่า \textit{doomscrolling} และมีรายงานว่าสัมพันธ์กับความวิตกกังวลที่สูงขึ้น
\cite{dominguez2025climate}

อย่างไรก็ตาม งานที่ใช้กลุ่มตัวอย่างขนาดใหญ่พบว่าความสัมพันธ์ระหว่างเวลาที่ใช้เทคโนโลยีดิจิทัล
กับสุขภาวะของวัยรุ่นมีขนาดเล็กมาก \cite{orben2019association}
และงานทบทวนวรรณกรรมสรุปว่าผลขึ้นกับรูปแบบการใช้งานและตัวผู้ใช้ ไม่ได้ขึ้นกับเวลาอย่างเดียว
\cite{sala2024umbrella}
ระหว่างเตรียมงาน ผู้วิจัยสำรวจชุดข้อมูลสาธารณะในหัวข้อนี้บนเว็บไซต์ Kaggle
\cite{singh_student_social_media,singh_social_media_addiction}
และพบว่าเป็นข้อมูลสังเคราะห์ที่สร้างไว้เพื่อการฝึกฝน ไม่ได้เก็บจากผู้ตอบจริง
จึงใช้อธิบายสถานการณ์ของนักศึกษาไม่ได้ ผู้วิจัยจึงเก็บข้อมูลเองจากนักศึกษาจริง

งานวิจัยนี้มีวัตถุประสงค์ 3 ข้อ ได้แก่
(1)~ศึกษาว่าเวลาที่ใช้โซเชียลมีเดียต่อวันสัมพันธ์กับภาวะซึมเศร้า ความวิตกกังวล และความเครียดหรือไม่
และมากน้อยเพียงใด
(2)~ศึกษาว่าประเภทเนื้อหาที่เสพ โดยเฉพาะการติดตามข่าวเชิงลบ สัมพันธ์กับสุขภาวะทางจิตแตกต่างกันหรือไม่
และ (3)~ศึกษาว่าความสัมพันธ์ดังกล่าวยังคงอยู่หรือไม่ เมื่อควบคุมการนอน การออกกำลังกาย
ความกดดันจากการเรียน ผลการเรียน และความเพียงพอทางการเงิน
งานวิจัยนี้ศึกษาเฉพาะความสัมพันธ์ ไม่ได้หาสาเหตุ และใช้กลุ่มตัวอย่างจากมหาวิทยาลัยเดียว
